\section{Evaluation and Benchmarks}
\label{sec:evaluation}

Evaluation must separate feedback \emph{quality}, whether a signal is accurate, calibrated, and unbiased, from feedback \emph{utilization}, whether the receiving system improves when that signal is supplied. The axes are independent: a correct critique can be ignored, while a reward model can score well on held-out preferences yet train a worse policy \citep{kim2024evaluating}. Treating end-task accuracy as a single summary obscures whether failure lies in the signal, the receiver, or their interface.

\subsection{Measuring Feedback Quality}
\label{subsec:eval-quality}

Quality benchmarks should match the signal's granularity and operational role. Outcome-preference benchmarks test reward-model ranking; RewardBench~2 is an anchor because it validates benchmark scores against best-of-$N$ selection and PPO outcomes rather than assuming pairwise accuracy transfers \citep{malik2025rewardbench}. Judge meta-evaluation tests verdict consistency and bias; JudgeBench ties labels to objective correctness and finds even strong judges only slightly above chance on hard pairs \citep{tan2024judgebench}. Critique benchmarks score error identification and correction; CriticBench evaluates generation, critique, and repair jointly \citep{lin2024criticbench}. Process benchmarks localize step-level errors; ProcessBench asks for the first erroneous reasoning step, directly testing the signal consumed by process supervision \citep{zheng2024processbench}. Embodied instruments test whether vision-language judges understand trajectories rather than final frames; ViSTa exposes severe weakness on sequential tasks \citep{wybitul2024vista}. These families are not interchangeable: a judge validated for reranking is not automatically a sound training reward or a useful critic.

Validity requires more than difficulty. Pairwise accuracy should predict the downstream use, and test distributions should expose length, style, refusal, and format. Process evaluation must separate error localization from final verdict because a wrong location misassigns credit. Generative critiques should separate factual correctness from actionable specificity. Report uncertainty and judge disagreement when small differences determine system selection.

\subsection{Measuring Feedback Utilization}
\label{subsec:eval-utilization}

Utilization benchmarks score the receiver under a controlled signal. RealCritic measures a critique through the correction it induces and shows that plausible open-loop critiques can fail to improve self-revision \citep{tang2025realcritic}. Interactive environments isolate feedback uptake across attempts; LLF-Bench uses post-action language as the learning signal and paraphrases feedback to prevent task recognition from masquerading as learning \citep{cheng2023llfbench}. Feedback Friction supplies near-oracle feedback yet finds that solvers still resist incorporating it, demonstrating that utilization can remain the bottleneck after quality is nearly solved \citep{jiang2025feedback}. Long-horizon settings add retrieval and temporal credit: MemoryAgentBench tests whether accumulated verbal experience is retrieved, used, and selectively forgotten across episodes \citep{hu2025evaluating}. The detailed benchmark matrix and evidence catalog are provided in the electronic supplement.

A clean utilization study varies the signal while holding task, receiver, and budget fixed. A no-feedback control measures prior capability; an oracle condition estimates attainable gain; irrelevant feedback tests resistance to persuasive noise. Report improvement per interaction or token, and test several receivers because one critique can help a capable model yet overwhelm a weaker one.

\subsection{Measurement Gaps and Minimum Reporting Protocol}
\label{subsec:eval-gaps}

Current evidence is strongest for short, verifiable, English, single-turn tasks and weakest for personalized, multilingual, long-horizon, adversarial, and embodied feedback. Judge-scored results inherit the judge's consistency, bias, and aggregation uncertainty. Embodied benchmarks largely grade signal quality and rarely test whether a motor policy can use verbal feedback. Static evaluation also misses failures that emerge only when outputs change later observations or rewards.

These gaps limit claims of generality. Results from a short answer verifier do not establish trajectory-level credit, and an English critique benchmark does not establish that the same feedback interface transfers across languages or cultures. Evaluation should therefore state the intended deployment regime and treat movement beyond the benchmark's signal type, horizon, modality, or interaction structure as an explicit extrapolation.

A minimum reporting protocol follows. (1)~Report feedback quality and utilization separately. (2)~Validate a quality benchmark against its downstream use, such as best-of-$N$ selection, policy optimization, correction success, or task completion. (3)~Audit judges and verifiers for consistency, bias, and adversarial robustness before using their verdicts as evaluation or reward. (4)~State the signal type, granularity, interaction stage, and receiver, so results transfer only to mechanisms with the same role. When possible, include a no-feedback control and an oracle-feedback upper bound; their gap separates absent information from failed uptake. This protocol does not yet solve the lack of standardized utilization metrics, but it prevents quality and utilization failures from being conflated.
